% ============================================================
% Paper draft for EMNLP 2026 System Demonstrations track.
%
% SETUP INSTRUCTIONS:
% 1. Download the official style files from:
%    https://github.com/acl-org/acl-style-files
%    (acl.sty, acl_natbib.bst, anthology.bib.txt -> anthology.bib)
% 2. Place this file (rename to e.g. paper.tex) in the same folder
%    as acl.sty.
% 3. Compile with pdflatex (recommended) + bibtex/biber.
% 4. Before final/camera-ready submission, change [review] to [final]
%    in the \usepackage line below (kept as [review] here since the
%    venue is single-blind and self-references are allowed, but the
%    "review" option controls page-numbering/anonymity formatting,
%    not blinding -- follow the venue's current instructions at
%    submission time).
% 5. Fill in REAL author names/affiliations/emails before submitting.
% 6. Fill in the actual license, live demo URL, and video URL in the
%    marked TODO locations before submitting -- these are STRICT,
%    ENFORCED requirements; the submission will be desk rejected
%    without them.
% ============================================================

\documentclass[11pt]{article}
\usepackage[preprint]{acl}
\usepackage{times}
\usepackage{latexsym}
\usepackage[T1]{fontenc}
\usepackage[utf8]{inputenc}
\usepackage{microtype}
\usepackage{inconsolata}
\usepackage{graphicx}
\usepackage{booktabs}

\title{A Multi-Signal Compliance Audit Pipeline for RAG Systems: \\
Combining NLI, Topical Similarity, and a Supervised Precedent Classifier}

% TODO: replace with real author names, affiliations, and emails.
% Reviewing is single-blind this year -- no anonymization needed.
\author{
Achyuta Raghunathan \\
UC Irvine\\
\texttt{achyutarag@gmail.com} \\
}

\begin{document}
\maketitle

\begin{abstract}
Compliance-answering systems built on retrieval-augmented generation (RAG)
are commonly audited with a single metric -- embedding cosine similarity or
natural language inference (NLI) argmax -- both of which have demonstrable
blind spots. Cosine similarity cannot detect a single inverted word (e.g.,
``waived'' vs.\ ``applied''), since the two sentences remain nearly identical
in vector space; NLI argmax collapses to ``compliant'' on any high-neutral
response, including evasive deflection and precedent-based policy
circumvention that never explicitly contradicts the source text. We present
a multi-signal audit pipeline that combines NLI entailment/contradiction
scoring, topical embedding similarity, and a small supervised classifier
targeting precedent/exception framing -- a failure mode we show neither NLI
nor similarity can resolve on its own. We report a fit/held-out/adversarial
evaluation protocol designed to prevent circular self-evaluation, and
document two discarded signal designs (a repurposed NLI probe and a
zero-shot classifier) alongside the working supervised approach, since the
negative results are informative about the failure modes of each technique.
The final system achieves 93.6\% accuracy (95\% CI [88.3\%, 98.7\%]) on
held-out and adversarial data, versus a 59.8\% naive NLI-argmax baseline
($p < 0.001$), with zero false positives and a documented, bounded set of
remaining false negatives.
\end{abstract}

\begin{center}
\small
\textbf{Links:} \; Code: \url{https://github.com/achyutarag/Multi-Signal-Compliance-Guard} \; $\vert$ \; Live Demo:
\url{https://multi-signal-compliance-guard-qq3xu59gdprwtww69gcrsh.streamlit.app/ } \; $\vert$ \; Video:
\url{https://www.loom.com/share/41b4122f9b3a46db84993814eded8a43}
\end{center}

\vspace{0.5em}
% NOTE: fill in the three URLs above as soon as each is available --
% repo can go in as soon as you push it (even before the demo/video are
% ready), so this line does not need to wait on the other two.

\section{Introduction}
\label{sec:intro}

\paragraph{Problem.} RAG and compliance-answering systems are commonly
audited with a single metric. In a regulated setting (banking, insurance,
healthcare), an approved answer that silently violates policy is a direct
compliance and liability exposure. We identify and evaluate against three
distinct failure modes that a single-metric audit misses: (1) direct
contradiction or inversion of policy content, (2) topic drift, in which a
response deflects onto an unrelated subject without confirming or denying
compliance, and (3) precedent framing, in which a response invokes
discretion, exception, or informal past practice to justify circumventing a
policy without directly contradicting its wording.

\paragraph{Importance and impact.} Compliance failures of type (2) and (3)
are, by construction, difficult for a human skimming a transcript to catch,
since the response never states an outright violation. An automated system
that can flag these cases for review has direct value for any organization
deploying RAG chatbots in a regulated domain.

\paragraph{Novelty.} Our contribution is not a single new model but a
\textbf{documented multi-signal decision procedure}, built through a
transparent, reproducible failure-mode taxonomy, together with a
fit/held-out/adversarial evaluation protocol that catches circular
self-evaluation before it reaches a significance claim. We report two
discarded signal designs (Sections~\ref{sec:v7}--\ref{sec:v8}) alongside the
working design (Section~\ref{sec:v9}), since we believe the negative results
are independently useful to practitioners considering similar approaches.

\paragraph{Target audience.} NLP engineers and applied researchers building
compliance/QA layers on top of RAG chatbots in regulated industries (finance,
insurance, healthcare, legal); researchers interested in evaluation
methodology for high-neutral NLI failure modes.

\section{Related Work / Comparison to Existing Approaches}
\label{sec:related}

\textbf{Cosine-similarity-only audits}, common in production RAG monitoring,
fail on negation and inversion, since semantically opposite statements about
the same entity can remain nearly identical in embedding space.

\textbf{NLI-argmax-only audits} fail on high-neutral evasive or
precedent-framing language: when a cross-encoder's neutral class dominates
the softmax output, an argmax decision rule maps the response to
``compliant'' by construction, regardless of whether the response is
evasive or discretion-invoking (Section~\ref{sec:v0}).

\textbf{Keyword/regex-based compliance filters}, common in legacy rule
engines, are shown directly in our evaluation (Section~\ref{sec:v3}) to be
actively harmful on adversarial phrasing: a hand-built list of trigger words
(``historically'', ``discretionary'', etc.) produced a $-68\%$ accuracy
delta versus a naive baseline on adversarial inputs specifically constructed
to contain those words while remaining policy-compliant.

\section{System Description}
\label{sec:system}

\subsection{Architecture}

The final system (v10) evaluates each (policy, response) pair through a
priority-ordered gate sequence, computing all signals unconditionally and
resolving the verdict via the first gate that fires:

\begin{enumerate}
\item \textbf{Direct contradiction gate.} If the NLI cross-encoder's
contradiction probability exceeds 0.35, the response is flagged as a
violation.
\item \textbf{Background-leak gate.} Within the high-neutral zone
(neutral probability $> 0.70$), if contradiction still exceeds 0.15, the
response is flagged.
\item \textbf{Precedent classifier gate.} A supervised logistic regression
classifier, trained on response embeddings and the (response $-$ policy)
embedding difference, flags responses whose predicted probability of
invoking precedent/exception framing exceeds 0.50.
\item \textbf{Topical similarity gate.} If the response's cosine similarity
to the policy text falls below a calibrated threshold (0.42), it is flagged
as topic drift.
\item \textbf{Direct margin fallback.} Outside the high-neutral zone, the
entailment$-$contradiction margin determines the verdict.
\end{enumerate}

We use \texttt{cross-encoder/nli-deberta-v3-base} \citep{he2021deberta,he2021debertav3} for NLI scoring and \texttt{all-MiniLM-L6-v2} \citep{wang2020minilm,reimers-gurevych-2019-sentence} for embeddings, both via the
\texttt{sentence-transformers} library \citep{reimers-gurevych-2019-sentence}, trained on the MultiNLI corpus \citep{williams-etal-2018-broad}.\footnote{Model cards:
\url{https://huggingface.co/cross-encoder/nli-deberta-v3-base} and
\url{https://huggingface.co/sentence-transformers/all-MiniLM-L6-v2}.}

\subsection{Why Two Signals Were Not Enough}
\label{sec:whyneed3rd}

An earlier version of the system (v6) used only the contradiction gate,
similarity gate, and margin fallback. In a larger subsequent evaluation
(Section~\ref{sec:v6limits}), this design caught only 19\% of precedent
framing violations. We traced this to overlapping similarity ranges: missed
violations had similarity scores of 0.432--0.675, directly overlapping with
genuinely compliant responses at 0.503--0.711. No similarity threshold can
separate these classes, because precedent-framing responses are, by
construction, topically on-topic with the policy -- the same property that
makes genuinely compliant paraphrases score highly. This motivated adding a
signal targeting the \emph{rhetorical structure} of the response rather than
its topical content (Sections~\ref{sec:v7}--\ref{sec:v9}).

\section{Evaluation Methodology}
\label{sec:eval-method}

\subsection{Fit / Held-out / Adversarial Protocol}

All thresholds and trained parameters are selected using only a \emph{fit}
split. Reported accuracy and significance figures are computed exclusively
on \emph{held-out} examples (same categories, unseen phrasing) and
\emph{adversarial} examples (constructed specifically to contain surface
features -- e.g., precedent-framing trigger words -- while having the
opposite ground-truth label). This separation was introduced after an
earlier design (Section~\ref{sec:v3}) achieved a spurious 100\% accuracy
(zero-width bootstrap confidence interval) by using a keyword list built
from reading the exact sentences it was later scored against.

\subsection{Dataset}

The evaluation dataset covers a single policy domain (a corporate
gift-value policy) across five categories -- Safe, Direct Failure, Topic
Drift, Precedent Framing, and Adversarial -- generated with LLM assistance
and verified example-by-example by a human annotator prior to any use in
threshold selection or training. Near-duplicate candidates (word-overlap
$> 0.55$) were identified and removed prior to verification. The final
dataset comprises 98 examples; further reduced to 93 after moving 5
adversarial examples into the classifier's training split
(Section~\ref{sec:v9}).

\subsection{Statistical Testing}

We report accuracy with 95\% confidence intervals from 1{,}000-iteration
paired bootstrap resampling \citep{berg-kirkpatrick-etal-2012-empirical},
and significance as the proportion of
bootstrap resamples in which the multi-signal system did not outperform the
naive baseline. We report per-category accuracy alongside all aggregate
figures, since an aggregate figure can mask a category-level failure when
categories are of unequal size and difficulty (Section~\ref{sec:v6limits}).

\section{Results}
\label{sec:results}

\subsection{v0: Initial Design (Invalid Result)}
\label{sec:v0}

An initial gate design required elevated contradiction probability to flag
drift, which genuinely evasive text does not exhibit (both entailment and
contradiction are low when neutral dominates). This caused the system to be
identical to the naive baseline on exactly the cases it was designed to
catch, producing a deterministic zero accuracy delta across all bootstrap
resamples ($p = 1.0000$) -- not evidence of no effect, but evidence of a
logic error.

\subsection{v3: Keyword Gate (Invalid Result)}
\label{sec:v3}

A subsequent design added a hand-built regular expression matching
precedent-framing trigger words. Evaluated only on the examples used to
build the word list, this achieved 100.00\% accuracy with a zero-width
bootstrap confidence interval ($p = 0.0000$) -- a signature of circular
evaluation rather than a genuine effect. Introducing held-out and
adversarial splits (Section~\ref{sec:eval-method}) revealed the true
behavior: the keyword gate produced a $-68.10\%$ accuracy delta versus
baseline on adversarial inputs specifically constructed to contain trigger
words in compliant contexts (e.g., ``we've traditionally reminded staff that
gifts over \$50 are never permitted'').

\subsection{v6: Similarity-Only Design, Limitations at Scale}
\label{sec:v6limits}

Table~\ref{tab:v6} shows the similarity-only design (contradiction gate +
similarity gate + margin fallback, no precedent-specific signal) evaluated
on a 98-example scaled dataset.

\begin{table}[t]
\centering
\small
\begin{tabular}{lrrr}
\toprule
\textbf{Category} & $n$ & Baseline & v6 Guard \\
\midrule
Safe               & 14 & 100.0\% & 100.0\% \\
Direct Failure     & 16 & 100.0\% & 100.0\% \\
Topic Drift        & 16 &   6.2\% &  93.8\% \\
Precedent Framing  & 16 &   0.0\% &  \textbf{18.8\%} \\
Adversarial        & 20 & 100.0\% &  70.0\% \\
\midrule
\textbf{Aggregate} & 82 &  62.3\% &  82.8\% \\
\bottomrule
\end{tabular}
\caption{Similarity-only design (v6/v9-precursor) on the scaled dataset.
Aggregate accuracy (82.8\%, $p<0.001$) masks near-total failure on
Precedent Framing (18.8\%), the category the system was designed to
address.}
\label{tab:v6}
\end{table}

This result illustrates a general risk in evaluation reporting: the
aggregate accuracy is statistically significant and numerically high, yet
the category central to the system's stated purpose is functioning barely
above the 0\% baseline. We report this explicitly rather than presenting
only the aggregate figure.

\subsection{v7: Repurposed NLI Probe (Discarded)}
\label{sec:v7}

We first attempted to address the precedent-framing gap by reusing the
existing NLI cross-encoder as a second probe, testing entailment between
the response and a synthetic claim describing the rhetorical move (``This
response describes an exception, discretionary carve-out, informal past
practice, or personal judgment call that deviates from the stated
policy\ldots''). On the fit split, the best achievable threshold accuracy
was 62.5\%, \emph{below} the 75\% majority-class baseline. Genuinely
compliant examples scored 0.93--0.99, indistinguishable from actual
violations. We attribute this to a distribution mismatch: NLI
cross-encoders are trained on concrete premise/hypothesis pairs describing
the same event or fact, not on judging abstract rhetorical categories.

\subsection{v8: Zero-Shot Classification (Discarded)}
\label{sec:v8}

We next tried Hugging Face's zero-shot-classification pipeline
(\texttt{facebook/bart-large-mnli} \citep{lewis-etal-2020-bart}), following the entailment-based zero-shot classification approach of \citet{yin-etal-2019-benchmarking}, with two mutually exclusive candidate labels. Fit-split accuracy exactly matched the 75\% majority-class
baseline, with scores compressed into a narrow, largely undiscriminating
band (0.659--0.996); compliant examples scored \emph{higher}
(0.976--0.990) than some actual violations (0.659). This is consistent
with known label-wording sensitivity in zero-shot MNLI classifiers.

\subsection{v9/v10: Supervised Precedent Classifier (Final Design)}
\label{sec:v9}

Both v7 and v8 share a root cause: relying on a pretrained model's
zero-shot judgment about an abstract category it was not trained to
evaluate. We instead trained a small, heavily regularized logistic
regression classifier \citep{pedregosa2011scikit} directly on fit-split examples, using response
embeddings and the (response $-$ policy) embedding difference as features.

An initial version (v9), trained only on Precedent Framing vs.\ Safe
examples, achieved 89.2\% aggregate accuracy (95\% CI [81.7\%, 95.1\%]) on
held-out and adversarial data ($n=82$), raising Precedent Framing accuracy
from 0\% to 87.5\%. However, it introduced 6 new false positives, all in
the Adversarial category, with scores of 0.501--0.510 -- directly on the
decision boundary. The classifier had never seen a single Adversarial-class
example during training, since that category was, by the original
evaluation design, routed entirely to the held-out adversarial split.

We corrected this (v10) by moving 5 Adversarial-category examples from the
adversarial split into the fit split and retraining on all three relevant
classes. Table~\ref{tab:final} reports the final result.

\begin{table}[t]
\centering
\small
\begin{tabular}{lrrr}
\toprule
\textbf{Category} & $n$ & Baseline & v10 Guard \\
\midrule
Safe               & 14 & 100.0\% & 100.0\% \\
Direct Failure     & 16 & 100.0\% & 100.0\% \\
Topic Drift        & 16 &   6.2\% &  93.8\% \\
Precedent Framing  & 16 &   0.0\% &  75.0\% \\
Adversarial        & 15 & 100.0\% & 100.0\% \\
\midrule
\textbf{Aggregate} & 77 &  59.8\% &  \textbf{93.6\%} \\
\bottomrule
\end{tabular}
\caption{Final system (v10) on held-out + adversarial data. 95\% CI
[88.3\%, 98.7\%] vs.\ baseline [49.4\%, 71.4\%], $p<0.001$. Zero false
positives; four remaining Precedent Framing false negatives, all within
0.03 of the decision boundary (0.467--0.498).}
\label{tab:final}
\end{table}

We note explicitly that the improvement from v9 to v10 involved a
disclosed trade-off: 2 correct Precedent Framing predictions were lost
(14/16 $\rightarrow$ 12/16) in exchange for eliminating all 6 Adversarial
false positives, and the Adversarial evaluation denominator shrank from 20
to 15. We report per-category counts rather than only the aggregate
percentage so this trade-off is independently verifiable.

\section{Licensing}

The evaluation code and pipeline are released under the MIT license. 
Underlying models are used under their respective Apache-2.0 open 
licenses: \texttt{cross-encoder/nli-deberta-v3-base} and 
\texttt{all-MiniLM-L6-v2} (see model cards, footnote 1).


\section*{Limitations}

\textbf{Single policy domain.} All evaluation examples derive from a single
gift-value policy. Generalization to other policy types (data handling,
expense approval, conflict-of-interest disclosure) is untested.

\textbf{Precedent-framing boundary calibration.} The four remaining false
negatives sit within 0.03 of the classifier's decision boundary
(0.467--0.498 vs.\ a 0.50 threshold), indicating a calibration gap rather
than a distinct unaddressed failure pattern. A larger, more diverse fit set
targeted at this boundary region is the most direct next step.

\textbf{No user study.} Evaluation is against hand-verified synthetic
examples, not real user or production traffic. Shadow-mode deployment
against live data, with human review retained in the loop, is recommended
before production use.

\textbf{Small fit sets.} Per-category fit sets are small (4--8 examples) by
design, to preserve held-out evaluation power. This is adequate for
threshold and boundary selection in our experiments, but means any single
mislabeled fit example has outsized influence; periodic re-auditing of
fit-set labels is recommended for any deployment.

\section*{Ethics Statement}

This system is designed to \emph{flag} candidate compliance violations for
human review, not to make unilateral compliance or disciplinary decisions.
Given the documented false-negative rate on precedent-framing language
(Table~\ref{tab:final}), we recommend deployments retain mandatory human
review for any flagged or borderline case, particularly in the 0.4--0.6
precedent-classifier score range where the model's confidence is lowest.
We used LLM-assisted generation for a portion of the evaluation dataset;
all generated examples were individually hand-verified for label accuracy
prior to any use in threshold selection, training, or reported evaluation,
to avoid circular or fabricated evaluation claims.

% ------------------------------------------------------------
% Add real \citep{}/\citet{} references as needed, e.g. for
% He et al. 2021 (DeBERTa), Reimers and Gurevych 2019
% (Sentence-BERT), Wang et al. 2020 (MiniLM), Lewis et al. 2020
% (BART). Populate custom.bib with verified BibTeX entries for
% any works you cite -- do not rely on this draft for exact
% bibliographic metadata.
% ------------------------------------------------------------
\bibliography{custom}

\appendix

\section{Dataset Generation and Hand-Verification Protocol}
\label{sec:appendix-gen}

Candidate examples were generated via LLM assistance, with category-specific
prompts targeting each of the five evaluation categories, in batches of 10
with a deduplication pass (word-overlap threshold 0.9) applied within and
across batches. All candidates were written to a review file containing the
proposed text, category, and proposed label, with \texttt{verified\_label},
\texttt{keep}, and \texttt{split} fields left blank. A human reviewer read
every candidate and populated these fields; the loading function used for
all downstream evaluation raises an error on any row missing a valid label
or split assignment, preventing an unreviewed file from silently entering
the evaluation pipeline.

\end{document}
